% !TEX program = xelatex
% !TeX encoding = UTF-8
\documentclass[12pt,a4paper]{article}
\usepackage{fontspec}
\usepackage{polyglossia}
\setmainlanguage{korean}
\setotherlanguage{english}

% ============================================================
% 한글 폰트 – 시스템에 설치된 Malgun Gothic 사용 (Windows)
% ============================================================
\setmainfont{Malgun Gothic}[Script=Hangul, Language=Korean]
\setsansfont{Malgun Gothic}[Script=Hangul, Language=Korean]
% 모노스페이스도 같은 한글 폰트로 (Consolas는 한글 미지원)
\setmonofont{Malgun Gothic}[Script=Hangul, Language=Korean]

% polyglossia에서 모노스페이스 한글을 위해 필요
\newfontfamily\hangulfonttt{Malgun Gothic}[Script=Hangul, Language=Korean]

% 영문 폰트 (수식 및 참고문헌 등)
\newfontfamily\englishfont{Times New Roman}[Ligatures=TeX]

% ============================================================
% 수학 및 서식 패키지
% ============================================================
\usepackage{amsmath,amssymb,amsthm,mathtools}
\usepackage{geometry}
\geometry{left=2.5cm,right=2.5cm,top=2.5cm,bottom=2.5cm}
\usepackage{booktabs}
\usepackage{hyperref}
\usepackage{url}
\usepackage{enumitem}
\usepackage{float}
\usepackage{xcolor}
\usepackage{titlesec}
\usepackage{fancyhdr}

% 색상 정의
\definecolor{gold}{RGB}{255,215,0}
\definecolor{darkblue}{RGB}{0,0,139}
\definecolor{skyblue}{RGB}{0,102,204}

% YUCT 매크로
\newcommand{\Keff}{K_{\mathrm{eff}}}
\newcommand{\kappac}{\kappa_c}
\newcommand{\eps}{\varepsilon}
\newcommand{\df}{d_{\!f}}
\newcommand{\dint}{\ensuremath{D\!+\!I\!\cdot\!R}}
\newcommand{\Sodd}{S_{\mathrm{odd}}}
\newcommand{\Seven}{S_{\mathrm{even}}}

% 한국어 정리 환경
\newtheorem{theorem}{정리}[section]
\newtheorem{definition}[theorem]{정의}
\newtheorem{proposition}[theorem]{명제}
\newtheorem{conjecture}[theorem]{추측}
\theoremstyle{remark}
\newtheorem{remark}[theorem]{비고}

\hypersetup{
	colorlinks=true,
	linkcolor=blue,
	citecolor=blue,
	urlcolor=blue,
}

% 제목 서식
\titleformat{\section}{\LARGE\bfseries\color{darkblue}}{\arabic{section}}{1em}{}
\titleformat{\subsection}{\Large\bfseries\color{darkblue}}{\arabic{subsection}}{1em}{}

% 바닥글
\fancypagestyle{firstpage}{%
	\fancyhf{}%
	\renewcommand{\headrulewidth}{0pt}%
	\renewcommand{\footrulewidth}{0pt}%
	\fancyfoot[C]{%
		\begin{minipage}[t]{\textwidth}
			\centering
			\textcolor{gold}{\rule{\textwidth}{1.6pt}}\\[5pt]
			{\small \textsuperscript{1}Yakushev Research, \url{https://yuct.org/}} \hfill
			{\small YPSDC 프로토콜: \url{https://ypsdc.com/}}\\[-2pt]
			\textcolor{gold}{\rule{\textwidth}{1.6pt}}\\[5pt]
			\textbf{Yakushev 통일 조정 이론 2026} \hfill \thepage\\[10pt]
		\end{minipage}%
	}%
}

\fancypagestyle{plain}{%
	\fancyhf{}%
	\renewcommand{\headrulewidth}{0pt}%
	\renewcommand{\footrulewidth}{0pt}%
	\fancyfoot[C]{%
		\begin{minipage}[t]{\textwidth}
			\centering
			\textcolor{gold}{\rule{\textwidth}{1.6pt}}\\[4pt]
			\textbf{Yakushev 통일 조정 이론 2026} \hfill \thepage\\[4pt]
		\end{minipage}%
	}%
}

\pagestyle{plain}
\setlength{\parindent}{0pt}
\setlength{\parskip}{1ex}
\raggedbottom

\makeatletter
\let\@ensure@LTR\relax
\makeatother

% ============================================================
% 문서 시작
% ============================================================
\begin{document}
	
	\begin{titlepage}
		\centering
		\vspace*{1cm}
		
		\Huge\textbf{부록 F (개정판): 경제 사전의 프랙탈 동역학과 화폐의 조정 이론}
		
		\vspace{0.5cm}
		\Large\textbf{환율, 통화 지역, 자가조직적 경제 경계에 대한 이론적 틀}
		
		\vspace{1.5cm}
		
		\Large\textbf{알렉세이 V. 야쿠셰프\textsuperscript{1}}
		
		YUCT 이론: \url{https://yuct.org/}\\
		YPSDC 프로토콜: \url{https://ypsdc.com/}
		
		\vspace{2cm}
		\textbf{DOI:} \url{https://doi.org/10.5281/zenodo.18444598}
		
		\vspace{1cm}
		
		\large
		2026년 6월 \\ \large 버전 3.0 (이전 부록 F를 대체함)
		
		\vspace{2cm}
		
		\begin{minipage}{0.9\textwidth}
			\begin{abstract}
				\noindent
				본 부록은 경제 분야에 대한 YUCT 접근법의 근본적인 재검토를 제시합니다. 우리는 이전의 모든 실증적 피팅 시도(예: GDP 또는 시장 변동성을 보편 지수 \(\beta = 0.67\)에 맞추려는 시도)를 포기하고, 대신 조정 존재론에 기반한 최초 원리로부터의 이론적 틀을 개발합니다. 경제는 상호작용하는 \textbf{경제 사전}(기업, 부문, 국가)의 네트워크로 모델링되며, 각 사전은 조정 효율성 \(K_{\mathrm{eff}}\)으로 특성화됩니다. 우리는 임의의 두 사전 사이의 경계가 만델브로트 집합과 유사한 프랙탈 동역학을 보임을 보여줍니다: 자기유사성, 복잡성, 그리고 보편 오차 법칙 \(\varepsilon = \kappa_c \alpha K_{\mathrm{eff}}^{-\beta}\)에 의해 지배됩니다. 화폐는 사전 경계를 가로지르는 활성화를 규제하는 \textbf{조정 매개변수}로 정의됩니다; 따라서 환율은 조정 효율성의 비율에 의해 결정됩니다. 이 틀은 몇 가지 정량적이고 반증 가능한 예측을 제공합니다: (1) 고주파 환율 시계열의 허스트 지수는 긴 시간 척도에서 \(H \approx 1/\beta \approx 1.5\)로 수렴해야 합니다; (2) 금융 수익률의 다중프랙탈 스펙트럼은 1차 대수 순환에 의해 결정되는 보편적 좌측 비대칭성을 보여야 합니다; (3) 환율의 변동성 \(\sigma\)는 두 사전 사이의 조정 간극 \(\Delta K_{\mathrm{eff}}\)에 대해 \(\sigma \propto (\Delta K_{\mathrm{eff}})^{-0.67}\)으로 스케일링됩니다. 본 부록에서는 실증적 피팅이 수행되지 않으며; 모든 예측은 경제 시스템에서 \(K_{\mathrm{eff}}\)에 대한 신뢰할 수 있는 프록시가 제공되면 테스트할 수 있는 형태로 공식화됩니다. 이 작업은 YUCT 경제학을 의심스러운 상관관계의 모음에서 엄격하고 반증 가능한 연구 프로그램으로 전환시킵니다.
			\end{abstract}
		\end{minipage}
		
		\vspace{1cm}
		
		\noindent
		\small\textbf{키워드:} YUCT, 조정 효율성, 경제 사전, 환율, 프랙탈 경계, 만델브로트 집합, 허스트 지수, 다중프랙탈 분석, 조정 매개변수로서의 화폐.
		
		\vspace{0.5cm}
		
		\noindent
		\textcopyright\ 2026 Yakushev Research. 모든 권리 보유.
	\end{titlepage}
	
	\tableofcontents
	\newpage
	
	% ============================================================
	% 1장: 서론 및 동기
	% ============================================================
	\section{서론: 새로운 시작이 필요한 이유}
	
	이 부록의 이전 버전들은 거시경제 시계열(예: GDP 성장 변동성, 암호화폐 예측 오차, 지정학적 위험 지표)이 YUCT의 보편 스케일링 법칙 \(\varepsilon \propto K_{\mathrm{eff}}^{-0.67}\)을 따른다는 것을 입증하려 했습니다. 이후의 독립적 검증 결과, 언급된 많은 표와 회귀분석이 오류가 있거나 사후에 유리한 사례를 선택한 것에 기반한 것으로 드러났습니다. 본 개정판은 이러한 결함을 완전히 인정하고 이전의 모든 실증적 주장을 철회합니다.
	
	근본적인 오류는 YUCT 프레임워크 자체에 있지 않고, '경제 사전'이 실제로 무엇이며 그 조정 효율성을 어떻게 측정할 수 있는지에 대한 적절한 이론적 모델을 먼저 개발하지 않은 채 직접적인 실증적 피팅을 강제하려 한 데 있었습니다. 경제학은 입자 물리학이 아닙니다: 경제 행위자는 점입자가 아니며, 그들의 질량은 기본 상수가 아닙니다. 따라서 경제 데이터에서 단순한 질량 척도를 찾는 것은 처음부터 오도된 것이었습니다.
	
	본 부록은 근본적으로 다른 접근법을 취합니다. 경제를 크기가 이산 스펙트럼에 속해야 하는 객체들의 집합으로 취급하는 대신, 우리는 그것을 상호작용하는 \textbf{사전들의 동적 네트워크}로 간주하며, 그 경계가 흥미로운 현상의 주요 장소입니다. 환율, 금리, 금융 변동성은 고립된 행위자의 속성이 아니라 사전 간 \textbf{경계의 창발적 속성}입니다. 우리는 이 경계들의 기하학이 프랙탈이고 자기유사하며, YUCT의 다른 모든 영역에 나타나는 동일한 보편 지수 \(\beta = 0.67\)에 의해 지배된다고 주장합니다.
	
	본 부록에서는 실증적 피팅이 전혀 수행되지 않습니다. 모든 정량적 주장은 YUCT 공리로부터 수학적으로 유도되며, 향후 연구를 위한 반증 가능한 예측으로 제시됩니다.
	
	% ============================================================
	% 2장: 경제 사전과 그 경계
	% ============================================================
	\section{경제 사전과 그 경계}
	
	\subsection{경제 사전의 정의}
	
	\begin{definition}[경제 사전]
		\textbf{경제 사전} \(\mathcal{D}_i\)은 경제 행위자(기업, 가구, 부문, 국가)가 외부 지표에 반응하여 활성화할 수 있는 가능한 행동, 계약, 생산 프로토콜, 가격 규칙 및 제도적 규범의 구조화된 집합입니다. 사전의 크기는 섀넌 엔트로피 \(H(\mathcal{D}_i)\)로 측정됩니다.
	\end{definition}
	
	예시:
	\begin{itemize}
		\item 기업의 사전에는 제품 카탈로그, 공급 계약, 가격 책정 알고리즘 및 내부 관리 절차가 포함됩니다.
		\item 국가 경제의 사전에는 법률 체계, 세법, 무역 협정, 통화 정책 프레임워크 및 금융 인프라가 포함됩니다.
	\end{itemize}
	
	사전 \(i\)의 조정 효율성은
	\begin{equation}
		K_{\mathrm{eff}}^{(i)} = \frac{H(\mathcal{A}_i)}{H(\kappa_i)},
		\label{eq:keff_econ}
	\end{equation}
	여기서 \(\mathcal{A}_i\)는 사전이 생성할 수 있는 행동의 집합이고, \(\kappa_i\)는 이를 유발하는 지표(가격 신호, 규제 지시, 시장 주문)입니다.
	
	\subsection{두 사전 사이의 경계}
	
	\begin{definition}[사전 경계]
		두 경제 사전 \(\mathcal{D}_i\)와 \(\mathcal{D}_j\) 사이의 \textbf{경계} \(\partial\mathcal{D}_{ij}\)는 한 사전에서 발생한 지표가 다른 사전의 입력을 활성화하는 모든 거래, 계약 및 정보 교환의 집합입니다.
	\end{definition}
	
	경계는 다음과 같은 현상의 자연스러운 장소입니다:
	\begin{itemize}
		\item 환율(국가 사전의 경우),
		\item 거래 비용 및 가격 차이(기업 수준 사전의 경우),
		\item 정보 비대칭 및 역선택.
	\end{itemize}
	
	\begin{proposition}[경계에서의 조정 매개변수]
		경계 \(\partial\mathcal{D}_{ij}\)를 통한 조정 효율성은 무차원 \textbf{조정 매개변수} \(\chi_{ij}\)로 측정되며, 이는 다음과 같이 정의됩니다:
		\begin{equation}
			\chi_{ij} = \frac{K_{\mathrm{eff}}^{(j)}}{K_{\mathrm{eff}}^{(i)}} \cdot \rho_{ij},
			\label{eq:chi}
		\end{equation}
		여기서 \(\rho_{ij}\)는 두 사전의 호환성(공통 표준, 공통 언어, 법적 조화 등)을 결정하는 공명 인자입니다.
	\end{proposition}
	
	\textbf{조정 매개변수로서의 화폐.} 독립 통화를 가진 두 국가 경제의 경우, 조정 매개변수 \(\chi_{ij}\)는 정확히 \textbf{환율}(통화 \(i\)당 통화 \(j\)의 단위)이며, 정규화 상수까지 같습니다. \(K_{\mathrm{eff}}^{(i)}\)에 비해 더 큰 \(K_{\mathrm{eff}}^{(j)}\)는 사전 \(j\)가 더 '가치 있음'을 의미하며 — 그 입력이 더 신뢰할 수 있게 활성화되므로 — 통화가 프리미엄을 요구합니다.
	
	% ============================================================
	% 3장: 경계의 프랙탈 동역학
	% ============================================================
	\section{경계의 프랙탈 동역학}
	
	\subsection{만델브로트 집합과의 유비}
	
	만델브로트 집합 \(\mathcal{M}\)은 반복 사상 \(z_{n+1} = z_n^2 + c\) (\(z_0 = 0\))으로 정의됩니다. 수열이 유계로 남아 있는 점 \(c\)는 \(\mathcal{M}\)에 속하고, 발산하는 점은 그 밖에 있습니다. 경계 \(\partial\mathcal{M}\)는 프랙탈이며, 무한히 복잡하고 모든 스케일에서 자기유사성을 보입니다.
	
	우리는 경제적 유비를 제안합니다. 두 사전 \(\mathcal{D}_i\)와 \(\mathcal{D}_j\)가 이산 시간 단계 \(n\)으로 인덱싱된 일련의 거래를 통해 상호작용한다고 합시다. \(n\)번의 거래 후 조정 상태는 복소 변수 \(z_n \in \mathbb{C}\)로 설명되며, 그 실수부와 허수부는 두 직교 차원(예: 가격 동기화 및 계약 이행)에서의 유사성 정도를 나타냅니다. 업데이트 규칙은
	\begin{equation}
		z_{n+1} = z_n^2 + \chi_{ij},
		\label{eq:mandel_econ}
	\end{equation}
	여기서 \(\chi_{ij} \in \mathbb{C}\)는 식~(\ref{eq:chi})의 (복소) 조정 매개변수입니다. 크기 \(|\chi_{ij}|\)는 결합 강도를 측정하고, 위상은 두 사전의 활성화 주기 사이의 위상 변이를 측정합니다.
	
	\begin{conjecture}[만델브로트의 경제적 대응]
		사전 쌍 \((\mathcal{D}_i, \mathcal{D}_j)\)은 \(\chi_{ij}\)가 일반화된 만델브로트 집합 \(\mathcal{M}_{\mathrm{econ}}\) 내에 있는 경우에만 유계이고 안정적인 조정을 유지할 수 있습니다(즉, 수열 \(\{z_n\}\)이 유계로 남음). 경계 \(\partial\mathcal{M}_{\mathrm{econ}}\)는 안정적인 조정 영역(안정적인 환율, 예측 가능한 무역 흐름)과 혼돈적 발산 영역(통화 위기, 초인플레이션, 무역 붕괴)을 구분합니다. 이 경계의 기하학은 프랙탈이며, 모든 시간 스케일에서 자기유사 구조를 가집니다.
	\end{conjecture}
	
	\subsection{경계에서의 보편 스케일링}
	
	경계에서의 동역학은 YUCT 전반에 적용되는 동일한 보편 조정 오차 법칙을 따르기 때문에, 경계 근처의 환율 통계적 특성은 보편 스케일링 법칙을 따라야 합니다. \(\sigma_{ij}(\Delta t)\)를 시간 지평 \(\Delta t\)에 걸쳐 측정된 사전 \(i\)와 \(j\) 사이의 환율 변동성이라고 합시다. 그러면 보편 오차 법칙 \(\varepsilon \propto K_{\mathrm{eff}}^{-\beta}\)로부터 다음을 얻습니다:
	
	\begin{proposition}[환율 변동성의 스케일링]
		조정 매개변수 \(\chi_{ij}\)의 변동성은 조정 간극 \(\Delta K_{\mathrm{eff}} = |K_{\mathrm{eff}}^{(i)} - K_{\mathrm{eff}}^{(j)}|\)에 대해 다음과 같이 스케일링됩니다:
		\begin{equation}
			\sigma(\chi_{ij}) \propto (\Delta K_{\mathrm{eff}})^{-\beta} = (\Delta K_{\mathrm{eff}})^{-0.67}.
			\label{eq:vol_scaling}
		\end{equation}
		매우 유사한 조정 효율성을 가진 사전 쌍(작은 간극 \(\Delta K_{\mathrm{eff}}\))은 높고 혼돈적인 변동성을 보이고; 큰 간극을 가진 쌍은 낮고 안정적인 변동성을 보입니다.
	\end{proposition}
	
	이것은 반증 가능한 예측입니다. \(K_{\mathrm{eff}}^{(i)}\)에 대한 신뢰할 수 있는 프록시(예: 세계은행의 거버넌스 지표와 무역 복잡성 데이터 결합)가 개발되면, 방정식~(\ref{eq:vol_scaling})은 국가 간 패널 데이터에 대해 테스트될 수 있습니다.
	
	\subsection{금융 시계열의 프랙탈 특성}
	경제물리학에서 잘 알려진 경험적 사실은 금융 시계열이 다중프랙탈 특성(변동성 군집, 장기 기억, 두꺼운 꼬리 수익률 분포)을 보인다는 것입니다. YUCT 내에서 이들은 우연이 아니라 사전 경계의 프랙탈 동역학의 필연적 결과입니다.
	
	\begin{proposition}[경계에서의 허스트 지수]
		안정적인 조정의 임계 경계 근처에서 작동하는 사전 쌍에 대해, 환율 시계열의 허스트 지수 \(H\)는
		\begin{equation}
			H \to \frac{1}{\beta} = \frac{3}{2} \quad \text{as} \quad \Delta t \to \infty.
			\label{eq:hurst}
		\end{equation}
		더 짧은 시간 스케일에서는 \(H\)가 더 낮은 값으로 교차하는 것을 보이며, 이는 경계의 다중프랙탈 특성을 반영합니다.
	\end{proposition}
	
	이 예측은 고주파 환율 데이터에 대한 표준 R/S 분석 또는 DFA로 테스트할 수 있습니다. \(H = 1.5\) 값은 매우 긴 기억을 가진 과정의 특징이며, 특이점에 가깝습니다 — 이는 대규모 금융 위기(예: 2008년 위기) 직전에 정확히 관찰되는 행동입니다.
	
	\subsection{단일 사전의 내부 프랙탈 조직}
	
	사전 간 경계가 프랙탈하고 혼돈적인 반면, 잘 조정된 사전의 내부 구조는 이산적인 계층적 질서를 보입니다. 사전의 항목들은 균일하게 분포하지 않고, 규칙, 계약 및 프로토콜의 자기유사 분기 트리를 형성합니다. 기업 규모, 소득 분포 및 도시 시스템의 계층 구조는 모두 이 내부 프랙탈 구조의 표현입니다.
	
	이전의 실패한 시도(전체 경제를 질량 척도에 맞추려는)와 달리, 우리는 이제 \textbf{하나의} 사전 내에서 하위 단위(기업, 도시, 소득 계층)의 크기가 프랙탈 차원 \(d_f\) 및 보편적 \(\beta\)와 관련된 지수를 가진 거듭제곱 법칙을 따라야 한다고 예측합니다. 도시에 대한 잘 알려진 지프 법칙(지수 \(\approx 1\))과 소득에 대한 파레토 법칙(지수 \(\approx 1.5\))은 고립된 기이함이 아니라 동일한 기본 조정 기하학의 서로 다른 투영입니다.
	
	% ============================================================
	% 4장: 수학적 틀
	% ============================================================
	\section{수학적 틀: 경계에서의 조정 라그랑지안}
	
	위에서 개괄한 정성적 그림을 수학적으로 정밀하게 만들기 위해, 우리는 이를 19차원 YUCT 라그랑주 형식론에 내장합니다. 경제 부문을 장 \(\Psi_{72}(x^\mu, y^m)\)로 기술합시다. 여기서 \(x^\mu\)는 4개의 시공간 좌표이고, \(y^m\)은 15개의 내부(콤팩트) 좌표입니다. 사전 \(\mathcal{D}_i\)는 특정 진공 기댓값 \(\langle \Psi_{72} \rangle_i\)에 해당하며, 두 사전 사이의 경계는 \(\langle \Psi_{72} \rangle_i\)와 \(\langle \Psi_{72} \rangle_j\) 사이를 보간하는 영역 벽 해입니다.
	
	경계에서의 유효 작용은
	\begin{equation}
		S_{\mathrm{boundary}} = \int d^4x \sqrt{-g} \left[ \frac{1}{2} \kappa (\nabla \chi)^2 + V(\chi) \right],
		\label{eq:action}
	\end{equation}
	여기서 \(\chi(x)\)는 국소 조정 매개변수(환율 장)이고, \(\kappa\)는 경계 강성이며, 퍼텐셜 \(V(\chi)\)는 다음과 같습니다:
	\begin{equation}
		V(\chi) = V_0 \left( e^{\lambda(\chi - \chi_0)} - \lambda(\chi - \chi_0) - 1 \right), \qquad \lambda = \frac{1}{\beta} = \frac{3}{2}.
		\label{eq:potential}
	\end{equation}
	이것은 중력 섹터(부록 G)에서 스칼라 조정 장을 지배하는 것과 동일한 퍼텐셜로, 보편성을 보장합니다.
	
	\(\chi\)에 대한 고전적 운동 방정식은
	\begin{equation}
		\kappa \Box \chi - V'(\chi) = 0,
		\label{eq:eom}
	\end{equation}
	이고, 그의 통계적 변동은
	\begin{equation}
		\langle (\delta \chi)^2 \rangle \propto \int \frac{d^4k}{(2\pi)^4} \frac{1}{\kappa k^2 + V''(\chi_0)} \propto \chi_0^{-2\beta},
		\label{eq:fluctuations}
	\end{equation}
	을 만족하며, 이는 식~(\ref{eq:vol_scaling})의 변동성 스케일링 법칙을 재생산합니다.
	
	% ============================================================
	% 5장: 반증 가능한 예측
	% ============================================================
	\section{반증 가능한 예측}
	
	위에서 개발된 틀은 데이터 피팅 없이 테스트할 수 있는 몇 가지 정량적 예측을 제공합니다:
	
	\begin{enumerate}[label=(\arabic*)]
		\item \textbf{환율의 허스트 지수.} 주요 자유 변동 통화쌍(예: EUR/USD, USD/JPY)에 대해, 20년간 일일 데이터로 계산된 장기 허스트 지수는 \(1.3 \le H \le 1.7\) 범위에 있어야 합니다. 주류 금융의 귀무 가설(랜덤 워크, \(H = 0.5\))은 높은 신뢰도로 기각되어야 합니다.
		
		\item \textbf{변동성-조정 간극 스케일링.} \(K_{\mathrm{eff}}\)에 대한 신뢰할 수 있는 프록시(예: 세계은행의 거버넌스 지표와 무역 복잡성 데이터 결합)가 구축되면, 환율의 횡단면 변동성은 절대 차이 \(\Delta K_{\mathrm{eff}}\)에 대해 \(\sigma \propto (\Delta K_{\mathrm{eff}})^{-0.67}\)로 스케일링되어야 합니다. 지수는 로그 좌표에서 선형 회귀로 추정할 수 있습니다; 예측은 기울기가 \(-0.67 \pm 0.10\)이라는 것입니다.
		
		\item \textbf{다중프랙탈 스펙트럼의 비대칭성.} 환율 수익률의 특이점 스펙트럼 \(f(\alpha)\)는 좌측 비대칭성을 보여야 하며, 비대칭 계수 \(A_\alpha \approx 0.30 \pm 0.05\)를 가져야 합니다. 이는 큰 조정 하락이 큰 조정 상승보다 우세함을 반영합니다. 이는 공개 데이터에 대한 표준 MF-DFA 프로그램으로 테스트할 수 있습니다.
		
		\item \textbf{통화 위기 전의 임계적 둔화.} 통화쌍이 임계 경계에 접근할 때(즉, \(\Delta K_{\mathrm{eff}} \to 0\)), 변동성의 자기상관 시간은 발산해야 합니다. 이는 주요 통화 위기 전에 허스트 지수의 꾸준한 증가와 다중프랙탈 스펙트럼의 수축이 관찰되어야 함을 예측하며, 이는 조기 경고 신호를 제공합니다.
	\end{enumerate}
	
	\textbf{예측의 상태:} 모든 예측은 경제 데이터에 대한 피팅 없이 YUCT 공리로부터 수학적으로 유도됩니다. 이들은 독립 연구자들이 직접 테스트할 수 있는 형태로 제시됩니다. 이들 중 하나라도 확인되면 화폐의 조정 이론에 대한 강력한 증거를 제공할 것입니다; 네 가지 모두의 반증은 틀의 근본적 수정을 요구할 것입니다.
	
	% ============================================================
	% 6장: 열린 문제 및 향후 과제
	% ============================================================
	\section{열린 문제 및 향후 과제}
	
	\begin{enumerate}[label=(\roman*)]
		\item \textbf{국가 \(K_{\mathrm{eff}}\)에 대한 신뢰할 수 있는 프록시 구축.} 가장 시급한 과제는 공개 데이터로부터 국가 경제의 조정 효율성을 추정하는 방법을 개발하는 것입니다. 후보: 법률 엔트로피, 수출 바스켓 복잡성(이달고-하우스만 지수), 제도적 품질(세계은행 거버넌스 지표), 시장 미세구조의 정보 측정치. 소규모 참조 국가 그룹에 대해 보정된 복합 지수는 전체 패널로 외삽될 수 있습니다.
		
		\item \textbf{경제적 만델브로트 사상의 수치 시뮬레이션.} 사상 \(z_{n+1} = z_n^2 + \chi\)을 체계적으로 연구해야 합니다. 어떤 \(\chi\) 값에 대해 수열이 유계로 남는가? 경계의 프랙탈 차원은 무엇인가? 잡음(유한 온도)을 추가하면 구조에 어떤 영향을 미치는가? 이러한 질문은 간단한 파이썬 스크립트로 처리할 수 있으며, 실제 환율 동역학과의 예상치 못한 정량적 연결을 드러낼 수 있습니다.
		
		\item \textbf{확립된 경제물리학과의 연결.} 5절에 나열된 예측은 경제물리학의 기존 결과(금융 수익률의 다중프랙탈성, 변동성의 장기 기억, 거듭제곱 법칙 꼬리)와 상당히 겹칩니다. 이 문헌을 YUCT 렌즈로 해석한 체계적 검토는 새로운 데이터 수집 없이 즉각적인 틀 검증을 제공할 것입니다.
	\end{enumerate}
	
	% ============================================================
	% 7장: 결론
	% ============================================================
	\section{결론}
	
	우리는 YUCT의 경제학 접근법을 완전히 재구성했습니다. 새로운 틀은 경제적 총량을 질량 척도에 강제로 맞추려는 실패한 전략을 포기하고, 대신 경제 사전과 그 경계에 대한 최초 원리로부터의 이론을 발전시킵니다. 화폐는 사전 간 활성화를 규제하는 조정 매개변수로 정의되며, 환율은 YUCT의 다른 모든 곳에 나타나는 동일한 보편 지수 \(\beta = 0.67\)에 의해 지배되는 것으로 나타납니다. 결과적인 예측은 정량적이고 반증 가능하며, 사후 데이터 피팅과 무관합니다. 이는 YUCT 경제학을 부담에서 유망하고 엄격한 연구 프로그램으로 전환시킵니다.
	
	\section*{면책 조항}
	
	본 부록은 이론적 작업입니다. 어떠한 실증적 주장도 하지 않으며, 투자 또는 정책 권고로 의도되지 않았습니다. 모든 정량적 예측은 과학적 검토를 위해 제시되며, 향후 데이터에 의해 반증될 수 있습니다.
	
	\section*{데이터 가용성}
	
	이 이론적 연구를 위해 새로운 데이터가 생성되지 않았습니다. 언급된 모든 출처는 공개적으로 이용 가능합니다.
	
	\begin{thebibliography}{99}
		\bibitem{AppendixA} A.V. Yakushev, ``부록 A: 19차원 YUCT 라그랑주 형식론,'' in \emph{YUCT}, Zenodo (2026).
		\bibitem{AppendixL} A.V. Yakushev, ``부록 L: 프랙탈 조정 오차 측정,'' in \emph{YUCT}, Zenodo (2026).
		\bibitem{AppendixG} A.V. Yakushev, ``부록 G: 기하학적 장력으로서의 중력,'' in \emph{YUCT}, Zenodo (2026).
		\bibitem{AppendixX} A.V. Yakushev, ``부록 X: 섀넌 정보 이론의 일반화,'' in \emph{YUCT}, Zenodo (2026).
		\bibitem{Mandelbrot1982} B.B. Mandelbrot, \emph{자연의 프랙탈 기하학}, W.H. Freeman (1982).
		\bibitem{Pareto1896} V. Pareto, \emph{정치경제학 강의}, Rouge (1896).
		\bibitem{Zipf1949} G.K. Zipf, \emph{인간 행동과 최소 노력의 원리}, Addison-Wesley (1949).
		\bibitem{Hidalgo2009} C.A. Hidalgo, R. Hausmann, ``경제적 복잡성의 구성 요소,'' \emph{PNAS} \textbf{106}, 10570–10575 (2009).
		\bibitem{Kantelhardt2002} J.W. Kantelhardt 외, ``비정상 시계열의 다중프랙탈 디트렌드 변동 분석,'' \emph{Physica A} \textbf{316}, 87–114 (2002).
	\end{thebibliography}
	
\end{document}